% ---------------------------------------------------------------
\section{Nepovinná úloha: meranie v jazyku PHP}

\subsection*{Projekt a nástroj}

Na nepovinnú úlohu som napísal malý e-shop v PHP s ôsmimi triedami:

\begin{itemize}
  \item \texttt{Product}, \texttt{DigitalProduct}, \texttt{Ebook} -- produkty, \texttt{Ebook} dedí od \texttt{DigitalProduct} a ten od \texttt{Product},
  \item \texttt{Customer} -- zákazník a jeho typ (\texttt{regular}, \texttt{student}, \texttt{vip}),
  \item \texttt{Cart} -- košík a súčet cien,
  \item \texttt{Order} -- objednávka, \texttt{calculateDiscount()} počíta zľavu podľa sumy, promo kódu, typu zákazníka a e-kníh,
  \item \texttt{InvoicePrinter} -- výpis faktúry,
  \item \texttt{Shop} -- spustí program a používa ostatné triedy.
\end{itemize}

Metriky som meral nástrojom PhpMetrics 2.11.0, ktorý vygeneruje HTML report. Hodnoty som čítal z jeho stránok. PhpMetrics nepočíta Coupling Factor, RFC ani CBO. Namiesto CBO uvádzam Efferent coupling, teda počet tried, ktoré daná trieda používa.

\subsection*{Maintainability Index}

\begin{figure}[H]
    \centering
    \includegraphics[width=1\linewidth]{bonus_mi.png}
    \caption{Maintainability Index tried v zozname ClassRank na úvodnej stránke reportu}
    \label{fig:bonus-mi}
\end{figure}

Hodnoty MI sú na úvodnej stránke reportu v zozname ClassRank (obr.~\ref{fig:bonus-mi}). Pri každej triede sú dve čísla: MI s komentármi a MI bez komentárov. V grafe vľavo je každá trieda kruh, veľkosť zodpovedá cyklomatickej zložitosti a farba MI.

\begin{center}
\begin{tabular}{llrr}
\toprule
 & Trieda & MI & MI bez komentárov \\
\midrule
Najlepšia & \texttt{Customer} & 78,83 & 61,85 \\
Najhoršia & \texttt{Order} & 47,38 & 38,18 \\
\bottomrule
\end{tabular}
\end{center}

\begin{figure}[H]
    \centering
    \includegraphics[width=1\linewidth]{bonus_loc.png}
    \caption{Logické riadky kódu (LLOC) a Halsteadov objem tried na stránke Size \& volume}
    \label{fig:bonus-loc}
\end{figure}

\texttt{Order} je v grafe najväčší červený kruh. Na stránke Size \& volume (obr.~\ref{fig:bonus-loc}) má najviac logických riadkov (69) aj najvyšší Halsteadov objem (593,88). \texttt{Customer} má iba konštruktor, dva gettery a jeden riadok komentára (CLOC = 1), preto je jeho MI s komentármi výrazne vyššie ako bez nich.

Hodnotu MI za celý projekt report neuvádza.

\subsection*{McCabe Cyclomatic Complexity}

\begin{figure}[H]
    \centering
    \includegraphics[width=1\linewidth]{bonus_cc.png}
    \caption{Cyklomatická zložitosť na stránke Complexity \& defects}
    \label{fig:bonus-cc}
\end{figure}

Hodnoty sú na stránke Complexity \& defects (obr.~\ref{fig:bonus-cc}). Stĺpec Max method cycl. je najvyššia CC jednej metódy v triede, WMC je súčet CC všetkých metód triedy.

\begin{center}
\begin{tabular}{llr}
\toprule
 & Trieda & Max method cycl. \\
\midrule
Najlepšia & \texttt{Product}, \texttt{Ebook}, \texttt{Shop}, \texttt{Customer}, \texttt{DigitalProduct} & 1 \\
Najhoršia & \texttt{Order} (metóda \texttt{calculateDiscount()}) & 17 \\
\bottomrule
\end{tabular}
\end{center}

\texttt{Order} má aj najvyššie WMC = 19. Takmer celá hodnota pochádza z metódy \texttt{calculateDiscount()}, v ktorej je reťazec \texttt{if / elseif} podľa sumy, \texttt{switch} na promo kódy, podmienky na typ zákazníka a cyklus cez položky košíka.

Pre celý projekt uvádza stránka priemernú cyklomatickú zložitosť na triedu 3,25 a priemerné WMC na triedu 4,13.

\subsection*{Depth of Inheritance Tree}

HTML report PhpMetrics hodnotu DIT neuvádza ani pre triedy, ani pre projekt. V kóde je jeden reťazec dedičnosti s tromi úrovňami: \texttt{Ebook} $\rightarrow$ \texttt{DigitalProduct} $\rightarrow$ \texttt{Product}. Ostatných päť tried nededí od žiadnej triedy.

\subsection*{Efferent coupling (namiesto CBO)}

\begin{figure}[H]
    \centering
    \includegraphics[width=1\linewidth]{bonus_coupling.png}
    \caption{Afferent a efferent coupling na stránke Coupling}
    \label{fig:bonus-coupling}
\end{figure}

Hodnoty sú na stránke Coupling (obr.~\ref{fig:bonus-coupling}). Efferent coupling je počet tried, ktoré daná trieda používa, afferent coupling počet tried, ktoré používajú ju.

\begin{center}
\begin{tabular}{llr}
\toprule
 & Trieda & Efferent coupling \\
\midrule
Najlepšia & \texttt{Product}, \texttt{Customer} & 0 \\
Najhoršia & \texttt{Shop} & 7 \\
\bottomrule
\end{tabular}
\end{center}

\texttt{Shop} v metóde \texttt{run()} vytvára košík, produkty, zákazníkov, objednávky a tlačiareň faktúr. \texttt{Product} a \texttt{Customer} žiadnu inú triedu nepoužívajú, iba ich používajú ostatné. \texttt{Product} používajú až tri triedy, má najvyšší afferent coupling (3).

Hodnotu za celý projekt stránka Coupling neuvádza.

\subsection*{Porovnanie s Java projektom}

\begin{center}
\begin{tabular}{lll}
\toprule
Metrika & Java (GameOfLife) & PHP (e-shop) \\
\midrule
Najhoršie MI triedy & \texttt{PatternLoader} 39,28 & \texttt{Order} 47,38 \\
Najvyššia CC metódy & \texttt{PatternLoader.load()} 20 & \texttt{Order::calculateDiscount()} 17 \\
Najvyššia väzba & \texttt{GameOfLifeApp} CBO 8 & \texttt{Shop} efferent 7 \\
\bottomrule
\end{tabular}
\end{center}

V oboch projektoch vyšla najhoršie trieda s jednou dlhou metódou plnou podmienok a najviac väzieb má trieda, ktorá spúšťa program. MetricsTree meria viac metrík ako PhpMetrics: počíta aj DIT, Coupling Factor, CBO a RFC.

\end{document}
